\section{Proposed framework}

MouseBrainBench is a modular framework for auditing partial digital models of
the mouse brain. It is organised around evidence gates. Each dataset contributes
a specific type of evidence and each possible claim is accepted only when the
required gates are satisfied. This design follows a conservative principle: the
framework must be able to produce positive results, but it must also be able to
reject claims when the evidence is incomplete.

The architecture is composed of five main modules. The first module is focused
on data integration. It stores provenance information and builds reproducible
adapters for public resources. The second module represents anatomical and
functional entities using explicit schemas. The third module builds the
benchmark tasks and baselines. The fourth module evaluates evidence through
metrics and the Mechanistic Identifiability Score. Finally, the fifth module
freezes the claims that can be defended from the generated artifacts.

\subsection{Design principles}

The first principle is to reuse existing resources instead of replacing them.
MouseBrainBench does not implement a new neuronal simulator, a new atlas, or a
new biological network format. Allen, Sensorium, Dynamic Sensorium, MICRONS, and
external simulators are treated as resources that can be connected to the
framework.

The second principle is to separate evidence types. Predictive performance,
reliability, anatomical plausibility, directed identifiability,
structure-function association, perturbation behaviour, and computational cost
are independent evidence blocks. This is relevant because a positive result in
one block does not necessarily validate the others.

The third principle is to make negative evidence explicit. In many benchmarks,
negative results are discarded because they do not improve the final score. In
MouseBrainBench, a negative result is informative when it prevents a wrong
interpretation. Therefore, Allen Visual Behavior Neuropixels is intentionally
included as a reproducible but non-identifiable case.

The fourth principle is reproducibility. The framework writes machine-readable
artifacts, keeps configuration files, and generates a publication freeze. This
freeze records the claims that are allowed and the claims that remain blocked at
the time of manuscript preparation.

\subsection{Mechanistic Identifiability Score}

The Mechanistic Identifiability Score is the decision mechanism of the
framework. It is conjunctive by design, meaning that a strong result in one
dimension cannot compensate for the absence of essential evidence in another
dimension. The score evaluates whether the model or analysis satisfies the
minimum conditions required for a specific claim.

This mechanism has an important practical consequence. Predictive models can be
accepted as useful predictors even when they are rejected as mechanistic
explanations. Similarly, local structure-function results can be accepted as
local observational evidence while being rejected as causal or whole-brain
claims. This distinction is the main safeguard of the proposal.

\begin{table}[t]
\centering
\caption{Dataset roles in MouseBrainBench. Each dataset contributes a distinct evidence layer rather than a generic positive claim.}
\label{tab:dataset-roles}
\small
\begin{tabularx}{\textwidth}{l l X X}
\toprule
Dataset & Scale & Role & Claim status \\
\midrule
Synthetic MIS & 4 controlled cases & Score sanity check & Method validation \\
Allen VBN & 20--21 sessions & Negative mechanistic control & Reproducible, not identifiable \\
Sensorium static & 7 validation mice & Predictive/reliability case & Predictive, not causal \\
Dynamic Sensorium & 5 current + 5 OOD mice & Temporal predictive stress test & Predictive, underpowered \\
MICRONS CAVE & 2991 units, 6575 synapses & Main empirical case & Replicated local association \\
\bottomrule
\end{tabularx}
\end{table}

